% Figure S2 legend — Nature-style supplementary figure caption
% Companion to: scripts/supplementary_figures/figure_s2.py

\textbf{Figure S2 $|$ End-to-end design and validation of the synthetic single-cell RNA-seq
data generator.}

\textbf{(A)}~Schematic of the cluster-plus-hub factor model used to define the gene-gene
correlation structure.
Each gene (blue nodes) belongs to a cluster whose members share variance through a
dedicated cluster hub (orange nodes).
A randomly selected subset of cluster hubs is additionally linked to a single global hub
(red star), introducing long-range correlation between clusters.
Edge weights are drawn from a uniform distribution and are modulated by the
shared-variance parameter $\alpha$.

\textbf{(B)}~Construction of the positive-definite correlation matrix $R$ from the loading
matrix $A$.
$A$ encodes the factor structure: an identity block (right, scaled by $\sqrt{1-\alpha}$)
provides each gene's unique variance, while additional columns (one per cluster factor and
one for the global hub) encode shared variance.
The correlation matrix is obtained as $R = \tilde{C}/\sqrt{\mathrm{diag}(\tilde{C})\,
\mathrm{diag}(\tilde{C})^\top}$, where $\tilde{C} = AA^\top$.

\textbf{(C,~D)}~Representative $100\times 100$ sub-blocks of correlation matrices generated
with high ($\alpha = 0.9$, \textbf{C}) and low ($\alpha = 0.5$, \textbf{D})
shared-variance fractions.
Block-diagonal structure reflects the cluster topology; off-diagonal inter-cluster
correlations arise through the global hub.

\textbf{(E)}~Flow diagram of count generation via a Gaussian copula.
\textit{Left:} Correlated latent variables are drawn from a multivariate normal
distribution $\mathcal{N}(\mathbf{0},\,R)$.
\textit{Center-left:} The standard normal CDF $\Phi$ maps each latent value to a
uniform marginal, preserving the rank-correlation structure.
\textit{Center-right:} Gene counts are obtained by applying the negative-binomial quantile
function $\mathrm{NB}^{-1}$ to each uniform variate; the per-gene mean $\mu_i$ is drawn
from an inverse-Gamma distribution to reproduce the heavy-tailed expression distribution
observed in real scRNA-seq data.
\textit{Right:} Cell-level library-size variability is introduced by multiplicative
log-normal amplitudes, followed by expression-dependent dropout
($P(\text{dropout}) = e^{-\lambda\,x}$) to generate the sparse observed count matrix.

\textbf{(F,~G)}~Complementary cumulative distribution functions (CCDFs) of per-cell
(\textbf{F}) and per-gene (\textbf{G}) total counts on log-log axes, confirming that the
simulated data exhibit the heavy-tailed expression distributions characteristic of
real single-cell data.

\textbf{(H)}~Binary (zero/non-zero) projection of a representative $120 \times 120$
sub-block of the observed count matrix, illustrating the high sparsity produced by the
dropout model.
The inset reports the overall fraction of zero entries.

\medskip
\noindent
All panels in rows~1--2 are generated from matrices with $n = 2000$ genes,
Pareto shape parameter $s = 1.5$, and hub-connectivity probability
$p_\mathrm{hub} = 0.2$.
Rows~3--4 use a representative simulation with $n_\mathrm{cells} = n_\mathrm{genes} = 400$,
$\alpha = 0.9$, dropout rate $\lambda = 1$, and inverse-Gamma shape/scale of $2/0.1$.